% !TeX root=../main.tex
\chapter{توسعه روش پیشنهادی و مطالعات تکمیلی}
\label{chap:extensions}

\section{مقدمه}
این فصل رساله را از یک مقالهٔ گسترش‌یافته فراتر می‌برد و تعمیم‌ها و مطالعات تکمیلی حول هستهٔ بازسازی کور مشترک را معرفی می‌کند.

\section{بازسازی در شرایط عدم همزمانی}
حالتی بررسی می‌شود که مرز دقیق نماد یا فریم در دسترس نباشد و همزمان‌سازی جزئی از مسئلهٔ بازسازی گردد.

\section{تخمین طول فریم ناشناخته}
وقتی \(L\) یا \(N\) از پیش معلوم نیست، راهبردهای تخمین طول و اتصال آن به جست‌وجوی ساختاری مطرح می‌شود.

\section{تخمین نرخ کد}
شناسایی نرخ مؤثر و ارتباط آن با الگوی افزونگی/پانچینگ بررسی می‌گردد.

\section{بازسازی تحت puncturing ناشناخته}
الگوی پانچینگ به مجموعهٔ پارامترهای مجهول افزوده می‌شود و اثر آن بر قیود و پیچیدگی تحلیل می‌شود.

\section{بازسازی با soft information}
استفاده از اطلاعات نرم برای بهبود آزمون قیود و کاهش پذیرش کاذب مطالعه می‌شود.

\section{بازسازی با تعداد فریم محدود}
عملکرد روش در رژیم دادهٔ کم و راهبردهای پایدارسازی آمار قیود بررسی می‌گردد.

\section{بازسازی در شرایط کانال متغیر}
کانال ناایستا یا تغییر آهستهٔ SNR چالش‌هایی برای آستانه‌های ثابت ایجاد می‌کند که در اینجا تحلیل می‌شود.

\section{استفاده از روش‌های هوشمند برای هدایت جست‌وجو}
رهیافت‌های مکاشفه‌ای برای اولویت‌بندی فرضیه‌ها و کاهش بسط غیرضروری معرفی می‌شوند.

\section{مقایسه Best-first، Beam Search و روش‌های دیگر}
راهبردهای جایگزین پیمایش درخت از نظر دقت، حافظه و زمان مقایسه می‌گردند.

\section{امکان استفاده از یادگیری ماشین برای pruning}
امکان جایگزینی یا تکمیل قواعد هرس کلاسیک با مدل‌های یادگیری برای پیش‌بینی نامزدهای امیدوارکننده بررسی می‌شود.
